
%=====================================================================
% Madhu Vadhani Madasamy - Iteration 3 Contributions
%=====================================================================

%=====================================================================
\section{Small Worked Example: One Complete ALS Iteration}
%=====================================================================
\textit{Owner Madhu Vadhani Madasamy. Checked by Margarita Kholostova and Aarav Jain.}

\subsection{Problem Setup}

To illustrate one complete Alternating Least Squares (ALS)
iteration, we consider a small rating matrix containing three
users and three movies, with two missing ratings.

The rating matrix is

\[
R =
\begin{bmatrix}
5 & ? & 1\\
4 & 2 & ?\\
4 & 3 & 5
\end{bmatrix}.
\]

The latent factor dimension is set to $k=1$, and the
regularization parameter is $\lambda=1$.

We initialize all movie latent factors to one:

\[
V^{(0)} =
\begin{bmatrix}
1\\
1\\
1
\end{bmatrix}.
\]

Missing ratings are excluded from the least-squares
calculations rather than treated as zero.

\subsection{Step 1: Updating User Factors}

Keeping the movie factors fixed, each user factor is
updated by minimizing the regularized least-squares
objective.

For $k=1$, the update is

\[
u_u =
\frac{\displaystyle\sum_{i\in\Omega_u}r_{ui}v_i}
{\displaystyle\lambda+
\sum_{i\in\Omega_u}v_i^2}.
\]

For User 1, the observed ratings are 5 and 1:

\[
u_1 =
\frac{5(1)+1(1)}{1+1^2+1^2}
=\frac{6}{3}=2.
\]

For User 2, the observed ratings are 4 and 2:

\[
u_2 =
\frac{4(1)+2(1)}{1+1^2+1^2}
=\frac{6}{3}=2.
\]

For User 3, the observed ratings are 4, 3 and 5:

\[
u_3 =
\frac{4(1)+3(1)+5(1)}{1+1^2+1^2+1^2}
=\frac{12}{4}=3.
\]

Therefore, the updated user factor vector is

\[
U^{(1)} =
\begin{bmatrix}
2\\
2\\
3
\end{bmatrix}.
\]

\subsection{Step 2: Updating Movie Factors}

Next, we keep the newly updated user factors fixed
and update each movie factor.

The update is

\[
v_i =
\frac{\displaystyle\sum_{u\in\Omega_i}r_{ui}u_u}
{\displaystyle\lambda+
\sum_{u\in\Omega_i}u_u^2}.
\]

For Movie 1, the observed ratings are 5, 4 and 4:

\[
v_1 =
\frac{5(2)+4(2)+4(3)}
{1+2^2+2^2+3^2}
=\frac{30}{18}=\frac{5}{3}.
\]

For Movie 2, the observed ratings are 2 and 3:

\[
v_2 =
\frac{2(2)+3(3)}
{1+2^2+3^2}
=\frac{13}{14}.
\]

For Movie 3, the observed ratings are 1 and 5:

\[
v_3 =
\frac{1(2)+5(3)}
{1+2^2+3^2}
=\frac{17}{14}.
\]

Therefore, the updated movie factor vector is

\[
V^{(1)} =
\begin{bmatrix}
\frac{5}{3}\\
\frac{13}{14}\\
\frac{17}{14}
\end{bmatrix}.
\]

\subsection{Final Results}

After one complete ALS iteration, the resulting
latent factor matrices are

\[
U^{(1)} =
\begin{bmatrix}
2\\
2\\
3
\end{bmatrix},
\qquad
V^{(1)} =
\begin{bmatrix}
\frac{5}{3}\\
\frac{13}{14}\\
\frac{17}{14}
\end{bmatrix}.
\]

The predicted rating matrix is obtained from

\[
\widehat{R}=U^{(1)}(V^{(1)})^\top.
\]

Thus,

\[
\widehat{R} =
\begin{bmatrix}
\frac{10}{3} & \mathbf{\frac{13}{7}} & \frac{17}{7}\\
\frac{10}{3} & \frac{13}{7} & \mathbf{\frac{17}{7}}\\
5 & \frac{39}{14} & \frac{51}{14}
\end{bmatrix}.
\]

The two entries in bold are the guesses for the missing
ratings: user 1 is expected to give movie 2
$\frac{13}{7}\approx1.86$ stars, and user 2 is expected to
give movie 3 $\frac{17}{7}\approx2.43$ stars.

\subsection{The Objective After Each Half Step}

With
$J(U,V)=\sum_{(u,i)\in\Omega}(r_{ui}-u_uv_i)^2
+\lambda\left(\sum_u u_u^2+\sum_i v_i^2\right)$,
the objective after the user update is

\[
\begin{aligned}
J(U^{(1)},V^{(0)})
&=\underbrace{9+1+4+0+1+0+4}_{\text{squared errors}}
+\underbrace{(4+4+9)+(1+1+1)}_{\text{penalty}}\\
&=19+20=39,
\end{aligned}
\]

and after the movie update it is

\[
J(U^{(1)},V^{(1)})=\frac{212}{7}\approx30.29.
\]

So $J$ goes down, as Question 4 of Iteration 2 says it
must.

The example demonstrates how ALS alternates between
user and movie updates while using only observed ratings.
All calculations are performed using exact fractions
to avoid floating-point rounding errors.

The implementation is provided in
\texttt{src/small\_example.py}.
The numerical results are independently checked against
the normal equations using NumPy's linear system solver
in \texttt{src/verify\_small\_example.py}. As a second way,
the same script runs one step of the team's ALS in
\texttt{src/methods.py}, which solves each system by our own
Gaussian elimination; it gives the same $U^{(1)}$ and
$V^{(1)}$ with a largest difference of 0.
